\chapter{Mechanized Operational Semantics in Lean 4}
\label{chap:lean_semantics}

\section{The Quest for Formally Verified Metacomputation}
\label{sec:lean_semantics:intro}

While compiler verification has achieved milestone triumphs in systems such as CompCert and CakeML, supercompilers have historically existed outside the reach of machine-checked proofs. The complex interactions between homeomorphic embedding whistles, process tree knot tying, and global generalization have resisted formalization.

NumLang achieves an unprecedented theoretical milestone: the entire operational semantics and core metacomputation transformations are mechanized in the Lean~4 proof assistant \citep{demoura2021lean}. All proof files reside in the repository under \texttt{lean/Supercompiler/}.

\section{Formal Specification of MIR in Lean 4}
\label{sec:lean_semantics:spec}

The abstract machine state and intermediate representations are mechanized directly in \texttt{lean/Supercompiler/Semantics.lean}:

\begin{lstlisting}[language=Lean4, caption={Mechanized MIR State and Values in Lean 4.}]
namespace Supercompiler

abbrev Local := String
abbrev BasicBlockId := Nat

inductive Op where
  | add | sub | mul | div | mod
  | bit_and | bit_or | bit_xor
  | eq | lt | ne | le | gt | ge
  deriving Repr, DecidableEq

inductive Val where
  | intVal : Int -> Val
  | ptrVal : Nat -> Val
  | boolVal : Bool -> Val
  deriving Repr, DecidableEq

abbrev MirEnv := List (Local * Val)
abbrev MirHeap := List Val

structure MirState where
  pc : BasicBlockId
  stmtIdx : Nat
  env : MirEnv
  heap : MirHeap
\end{lstlisting}

\section{The Small-Step Transition Relation (Step)}
\label{sec:lean_semantics:step}

The small-step execution of SSA MIR is formalized as an inductive proposition $\mathtt{Step} : \mathtt{MirFunction} \to \mathtt{MirState} \to \mathtt{MirState} \to \mathtt{Prop}$:

\begin{lstlisting}[language=Lean4, caption={The Step relation in Lean 4.}]
inductive Step (fn : MirFunction) : MirState -> MirState -> Prop where
  | stmt (s : MirState) (b : MirBasicBlock) (st : Statement) (env' : MirEnv) :
      b \in fn.blocks ->
      b.id = s.pc ->
      getStmt b.stmts s.stmtIdx = some st ->
      StmtStep s.env st env' ->
      Step fn s { s with stmtIdx := s.stmtIdx + 1, env := env' }
  | branch (s : MirState) (b : MirBasicBlock) (target : BasicBlockId) :
      b \in fn.blocks ->
      b.id = s.pc ->
      s.stmtIdx = b.stmts.length ->
      b.term = Terminator.Branch target ->
      Step fn s { s with pc := target, stmtIdx := 0 }
  | branchIf_true (s : MirState) (b : MirBasicBlock) (cond : Local) (then_t else_t : BasicBlockId) (n : Int) :
      b \in fn.blocks ->
      b.id = s.pc ->
      s.stmtIdx = b.stmts.length ->
      b.term = Terminator.BranchIf cond then_t else_t ->
      lookup s.env cond = some (Val.intVal n) ->
      n != 0 ->
      Step fn s { s with pc := then_t, stmtIdx := 0 }
\end{lstlisting}

\section{The Constructive Multi-Step Relation (Evaluates)}
\label{sec:lean_semantics:evaluates}

In Phase 42, NumLang addressed a subtle circularity present in prior academic formalizations. Classical definitions often define evaluation via mutual induction with function calls, creating non-well-founded premises.

NumLang resolves this by defining evaluation constructively via the reflexive-transitive closure $\mathtt{StepStar}$:

\begin{lstlisting}[language=Lean4, caption={Evaluates and SemanticEquivalence in Lean 4.}]
inductive StepStar (fn : MirFunction) : MirState -> MirState -> Prop where
  | refl (s : MirState) : StepStar fn s s
  | step (s1 s2 s3 : MirState) : Step fn s1 s2 -> StepStar fn s2 s3 -> StepStar fn s1 s3

def Evaluates (fn : MirFunction) (args : MirEnv) (res : Val) : Prop :=
  \exists s_final, StepStar fn (initState fn args) s_final /\ TerminatesWith fn s_final res

def SemanticEquivalent (f1 f2 : MirFunction) : Prop :=
  \forall args res, Evaluates f1 args res <-> Evaluates f2 args res
\end{lstlisting}

This constructive formulation provides the unshakeable foundation for all downstream preservation proofs.

\section{Verbatim Mechanization: Supercompiler/Semantics.lean}
\label{sec:lean_semantics:verbatim}

The complete mechanized specification of MIR syntax, values, and operational semantics is reproduced below directly from \texttt{lean/Supercompiler/Semantics.lean}:

\begin{lstlisting}[language=Lean4, caption={Full Mechanized Semantics in Lean 4 (\texttt{lean/Supercompiler/Semantics.lean}).}]
namespace Supercompiler

abbrev Local := String
abbrev BasicBlockId := Nat

inductive Op where
  | add
  | sub
  | mul
  | div
  | mod
  | bit_and
  | bit_or
  | bit_xor
  | eq
  | lt
  | ne
  | le
  | gt
  | ge
  deriving Repr, DecidableEq

inductive Val where
  | intVal : Int -> Val
  | ptrVal : Nat -> Val
  | boolVal : Bool -> Val
  deriving Repr, DecidableEq

abbrev MirEnv := List (Local  x  Val)
abbrev MirHeap := List Val

def lookup (env : MirEnv) (x : Local) : Option Val :=
  match env with
  | [] => none
  | (y, v) :: rest => if x = y then some v else lookup rest x

def update (env : MirEnv) (x : Local) (v : Val) : MirEnv :=
  (x, v) :: env

def intToUInt64 (v : Int) : UInt64 :=
  let mod2_64 : Nat := 18446744073709551616
  if v >= 0 then
    UInt64.ofNat (v.toNat % mod2_64)
  else
    let rem := v.natAbs % mod2_64
    if rem = 0 then 0
    else UInt64.ofNat (mod2_64 - rem)

def uInt64ToInt (u : UInt64) : Int :=
  let n := u.toNat
  let pow2_63 : Nat := 9223372036854775808
  let pow2_64 : Nat := 18446744073709551616
  if n >= pow2_63 then
    - (Int.ofNat (pow2_64 - n))
  else
    Int.ofNat n

def evalOp (op : Op) (v1 v2 : Int) : Option Int :=
  match op with
  | Op.add => some (v1 + v2)
  | Op.sub => some (v1 - v2)
  | Op.mul => some (v1 * v2)
  | Op.div =>
    if v2 = 0 then none
    else
      let d := Int.ofNat (v1.natAbs / v2.natAbs)
      some (if (v1 < 0) == (v2 < 0) then d else -d)
  | Op.mod =>
    if v2 = 0 then none
    else
      let rem := Int.ofNat (v1.natAbs % v2.natAbs)
      some (if v1 < 0 then -rem else rem)
  | Op.bit_and => some (uInt64ToInt (intToUInt64 v1 &&& intToUInt64 v2))
  | Op.bit_or  => some (uInt64ToInt (intToUInt64 v1 ||| intToUInt64 v2))
  | Op.bit_xor => some (uInt64ToInt (intToUInt64 v1 ^^^ intToUInt64 v2))
  | Op.eq  => some (if v1 = v2 then 1 else 0)
  | Op.lt  => some (if v1 < v2 then 1 else 0)
  | Op.ne  => some (if v1 != v2 then 1 else 0)
  | Op.le  => some (if v1 <= v2 then 1 else 0)
  | Op.gt  => some (if v1 > v2 then 1 else 0)
  | Op.ge  => some (if v1 >= v2 then 1 else 0)

inductive Rvalue where
  | Use : Local -> Rvalue
  | Constant : Int -> Rvalue
  | BinOp : Op -> Local -> Local -> Rvalue
  | Call : String -> List Local -> Rvalue
  | Alloc : Local -> Rvalue
  | Load : Local -> Rvalue
  deriving Repr, DecidableEq

inductive Statement where
  | Assign : Local -> Rvalue -> Statement
  | Nop : Statement
  deriving Repr, DecidableEq

inductive Terminator where
  | Branch : BasicBlockId -> Terminator
  | BranchIf : Local -> BasicBlockId -> BasicBlockId -> Terminator
  | Switch : Local -> List (Int  x  BasicBlockId) -> BasicBlockId -> Terminator
  | Return : Option Local -> Terminator
  | Unreachable : Terminator
  | Fork : BasicBlockId -> BasicBlockId -> BasicBlockId -> Terminator
  deriving Repr, DecidableEq

structure MirBasicBlock where
  id : BasicBlockId
  stmts : List Statement
  term : Terminator
  deriving Repr, DecidableEq

structure MirFunction where
  name : String
  params : List Local
  blocks : List MirBasicBlock
  entry : BasicBlockId
  deriving Repr, DecidableEq

structure MirProgram where
  functions : List MirFunction
  deriving Repr, DecidableEq

structure MirState where
  pc : BasicBlockId
  stmtIdx : Nat
  env : MirEnv
  heap : MirHeap
  deriving Repr, DecidableEq

inductive StmtStep : MirEnv -> Statement -> MirEnv -> Prop where
  | assign_const (env : MirEnv) (x : Local) (n : Int) :
      StmtStep env (Statement.Assign x (Rvalue.Constant n)) (update env x (Val.intVal n))
  | assign_use (env : MirEnv) (x y : Local) (v : Val) :
      lookup env y = some v ->
      StmtStep env (Statement.Assign x (Rvalue.Use y)) (update env x v)
  | assign_binop (env : MirEnv) (x y z : Local) (op : Op) (n1 n2 n3 : Int) :
      lookup env y = some (Val.intVal n1) ->
      lookup env z = some (Val.intVal n2) ->
      evalOp op n1 n2 = some n3 ->
      StmtStep env (Statement.Assign x (Rvalue.BinOp op y z)) (update env x (Val.intVal n3))
  | assign_call (env : MirEnv) (x : Local) (f : String) (args : List Local) (retVal : Val) :
      StmtStep env (Statement.Assign x (Rvalue.Call f args)) (update env x retVal)
  | nop (env : MirEnv) :
      StmtStep env Statement.Nop env

def getStmt (stmts : List Statement) (idx : Nat) : Option Statement :=
  match stmts, idx with
  | [], _ => none
  | x :: _, 0 => some x
  | _ :: xs, n + 1 => getStmt xs n

inductive Step (fn : MirFunction) : MirState -> MirState -> Prop where
  | stmt (s : MirState) (b : MirBasicBlock) (st : Statement) (env' : MirEnv) :
      b  in  fn.blocks 
\end{lstlisting}

\section{Verbatim Operational Semantics in Lean 4 (Semantics.lean)}
\label{sec:lean_sem:verbatim_listing}

Listing~\ref{lst:semantics_lean_verbatim} contains the complete formalization of the NumLang operational semantics from \texttt{lean/Supercompiler/Semantics.lean}:

\begin{lstlisting}[language=Lean4, caption={Complete Formal Operational Semantics Mechanization in Lean 4 (\texttt{lean/Supercompiler/Semantics.lean}).}, label={lst:semantics_lean_verbatim}]
namespace Supercompiler

abbrev Local := String
abbrev BasicBlockId := Nat

inductive Op where
  | add
  | sub
  | mul
  | div
  | mod
  | bit_and
  | bit_or
  | bit_xor
  | eq
  | lt
  | ne
  | le
  | gt
  | ge
  deriving Repr, DecidableEq

inductive Val where
  | intVal : Int -> Val
  | ptrVal : Nat -> Val
  | boolVal : Bool -> Val
  deriving Repr, DecidableEq

abbrev MirEnv := List (Local  x  Val)
abbrev MirHeap := List Val

def lookup (env : MirEnv) (x : Local) : Option Val :=
  match env with
  | [] => none
  | (y, v) :: rest => if x = y then some v else lookup rest x

def update (env : MirEnv) (x : Local) (v : Val) : MirEnv :=
  (x, v) :: env

def intToUInt64 (v : Int) : UInt64 :=
  let mod2_64 : Nat := 18446744073709551616
  if v >= 0 then
    UInt64.ofNat (v.toNat % mod2_64)
  else
    let rem := v.natAbs % mod2_64
    if rem = 0 then 0
    else UInt64.ofNat (mod2_64 - rem)

def uInt64ToInt (u : UInt64) : Int :=
  let n := u.toNat
  let pow2_63 : Nat := 9223372036854775808
  let pow2_64 : Nat := 18446744073709551616
  if n >= pow2_63 then
    - (Int.ofNat (pow2_64 - n))
  else
    Int.ofNat n

def evalOp (op : Op) (v1 v2 : Int) : Option Int :=
  match op with
  | Op.add => some (v1 + v2)
  | Op.sub => some (v1 - v2)
  | Op.mul => some (v1 * v2)
  | Op.div =>
    if v2 = 0 then none
    else
      let d := Int.ofNat (v1.natAbs / v2.natAbs)
      some (if (v1 < 0) == (v2 < 0) then d else -d)
  | Op.mod =>
    if v2 = 0 then none
    else
      let rem := Int.ofNat (v1.natAbs % v2.natAbs)
      some (if v1 < 0 then -rem else rem)
  | Op.bit_and => some (uInt64ToInt (intToUInt64 v1 &&& intToUInt64 v2))
  | Op.bit_or  => some (uInt64ToInt (intToUInt64 v1 ||| intToUInt64 v2))
  | Op.bit_xor => some (uInt64ToInt (intToUInt64 v1 ^^^ intToUInt64 v2))
  | Op.eq  => some (if v1 = v2 then 1 else 0)
  | Op.lt  => some (if v1 < v2 then 1 else 0)
  | Op.ne  => some (if v1 != v2 then 1 else 0)
  | Op.le  => some (if v1 <= v2 then 1 else 0)
  | Op.gt  => some (if v1 > v2 then 1 else 0)
  | Op.ge  => some (if v1 >= v2 then 1 else 0)

inductive Rvalue where
  | Use : Local -> Rvalue
  | Constant : Int -> Rvalue
  | BinOp : Op -> Local -> Local -> Rvalue
  | Call : String -> List Local -> Rvalue
  | Alloc : Local -> Rvalue
  | Load : Local -> Rvalue
  deriving Repr, DecidableEq

inductive Statement where
  | Assign : Local -> Rvalue -> Statement
  | Nop : Statement
  deriving Repr, DecidableEq

inductive Terminator where
  | Branch : BasicBlockId -> Terminator
  | BranchIf : Local -> BasicBlockId -> BasicBlockId -> Terminator
  | Switch : Local -> List (Int  x  BasicBlockId) -> BasicBlockId -> Terminator
  | Return : Option Local -> Terminator
  | Unreachable : Terminator
  | Fork : BasicBlockId -> BasicBlockId -> BasicBlockId -> Terminator
  deriving Repr, DecidableEq

structure MirBasicBlock where
  id : BasicBlockId
  stmts : List Statement
  term : Terminator
  deriving Repr, DecidableEq

structure MirFunction where
  name : String
  params : List Local
  blocks : List MirBasicBlock
  entry : BasicBlockId
  deriving Repr, DecidableEq

structure MirProgram where
  functions : List MirFunction
  deriving Repr, DecidableEq

structure MirState where
  pc : BasicBlockId
  stmtIdx : Nat
  env : MirEnv
  heap : MirHeap
  deriving Repr, DecidableEq

inductive StmtStep : MirEnv -> Statement -> MirEnv -> Prop where
  | assign_const (env : MirEnv) (x : Local) (n : Int) :
      StmtStep env (Statement.Assign x (Rvalue.Constant n)) (update env x (Val.intVal n))
  | assign_use (env : MirEnv) (x y : Local) (v : Val) :
      lookup env y = some v ->
      StmtStep env (Statement.Assign x (Rvalue.Use y)) (update env x v)
  | assign_binop (env : MirEnv) (x y z : Local) (op : Op) (n1 n2 n3 : Int) :
      lookup env y = some (Val.intVal n1) ->
      lookup env z = some (Val.intVal n2) ->
      evalOp op n1 n2 = some n3 ->
      StmtStep env (Statement.Assign x (Rvalue.BinOp op y z)) (update env x (Val.intVal n3))
  | assign_call (env : MirEnv) (x : Local) (f : String) (args : List Local) (retVal : Val) :
      StmtStep env (Statement.Assign x (Rvalue.Call f args)) (update env x retVal)
  | nop (env : MirEnv) :
      StmtStep env Statement.Nop env

def getStmt (stmts : List Statement) (idx : Nat) : Option Statement :=
  match stmts, idx with
  | [], _ => none
  | x :: _, 0 => some x
  | _ :: xs, n + 1 => getStmt xs n

inductive Step (fn : MirFunction) : MirState -> MirState -> Prop where
  | stmt (s : MirState) (b : MirBasicBlock) (st : Statement) (env' : MirEnv) :
      b  in  fn.blocks ->
      b.id = s.pc ->
      getStmt b.stmts s.stmtIdx = some st ->
      StmtStep s.env st env' ->
      Step fn s { s with stmtIdx := s.stmtIdx + 1, env := env' }
  | branch (s : MirState) (b : MirBasicBlock) (target : BasicBlockId) :
      b  in  fn.blocks ->
      b.id = s.pc ->
      s.stmtIdx = b.stmts.length ->
      b.term = Terminator.Branch target ->
      Step fn s { s with pc := target, stmtIdx := 0 }
  | branchIf_true (s : MirState) (b : MirBasicBlock) (cond : Local) (then_t else_t : BasicBlockId) (n : Int) :
      b  in  fn.blocks ->
      b.id = s.pc ->
      s.stmtIdx = b.stmts.length ->
      b.term = Terminator.BranchIf cond then_t else_t ->
      lookup s.env cond = some (Val.intVal n) ->
      n != 0 ->
      Step fn s { s with pc := then_t, stmtIdx := 0 }
  | branchIf_false (s : MirState) (b : MirBasicBlock) (cond : Local) (then_t else_t : BasicBlockId) :
      b  in  fn.blocks ->
      b.id = s.pc ->
      s.stmtIdx = b.stmts.length ->
      b.term = Terminator.BranchIf cond then_t else_t ->
      lookup s.env cond = some (Val.intVal 0) ->
      Step fn s { s with pc := else_t, stmtIdx := 0 }
  | switch (s : MirState) (b : MirBasicBlock) (var : Local) (targets : List (Int  x  BasicBlockId)) (default_t : BasicBlockId) (target : BasicBlockId) :
      b  in  fn.blocks ->
      b.id = s.pc ->
      s.stmtIdx = b.stmts.length ->
      b.term = Terminator.Switch var targets default_t ->
      Step fn s { s with pc := target, stmtIdx := 0 }
  | fork (s : MirState) (b : MirBasicBlock) (left right join : BasicBlockId) :
      b  in  fn.blocks ->
      b.id = s.pc ->
      s.stmtIdx = b.stmts.length ->
      b.term = Terminator.Fork left right join ->
      Step fn s { s with pc := join, stmtIdx := 0 }

inductive TerminatesWith (fn : MirFunction) (s : MirState) (res : Val) : Prop where
  | ret_some (b : MirBasicBlock) (retVar : Local) :
      b  in  fn.blocks ->
      b.id = s.pc ->
      s.stmtIdx = b.stmts.length ->
      b.term = Terminator.Return (some retVar) ->
      lookup s.env retVar = some res ->
      TerminatesWith fn s res
  | ret_none (b : MirBasicBlock) :
      b  in  fn.blocks ->
      b.id = s.pc ->
      s.stmtIdx = b.stmts.length ->
      b.term = Terminator.Return none ->
      res = Val.intVal 0 ->
      TerminatesWith fn s res

inductive StepStar (fn : MirFunction) : MirState -> MirState -> Prop where
  | refl (s : MirState) : StepStar fn s s
  | step (s1 s2 s3 : MirState) : Step fn s1 s2 -> StepStar fn s2 s3 -> StepStar fn s1 s3

theorem stepstar_trans {fn : MirFunction} {s1 s2 s3 : MirState}
    (h1 : StepStar fn s1 s2) (h2 : StepStar fn s2 s3) :
    StepStar fn s1 s3 := by
  induction h1 with
  | refl _ => exact h2
  | step s_a s_b s_c hstep _ ih =>
    exact StepStar.step s_a s_b s3 hstep (ih h2)

theorem stepstar_single {fn : MirFunction} {s1 s2 : MirState}
    (h : Step fn s1 s2) : StepStar fn s1 s2 :=
  StepStar.step s1 s2 s2 h (StepStar.refl s2)

def initState (fn : MirFunction) (args : MirEnv) : MirState :=
  { pc := fn.entry, stmtIdx := 0, env := args, heap := [] }

def Evaluates (fn : MirFunction) (args : MirEnv) (res : Val) : Prop :=
  exists  s_final, StepStar fn (initState fn args) s_final /\ TerminatesWith fn s_final res

def SemanticEquivalent (f1 f2 : MirFunction) : Prop :=
  forall  args res, Evaluates f1 args res <-> Evaluates f2 args res

theorem semantic_equiv_refl (f : MirFunction) : SemanticEquivalent f f := by
  intro args res
  exact Iff.rfl

theorem semantic_equiv_symm {f1 f2 : MirFunction} (h : SemanticEquivalent f1 f2) :
    SemanticEquivalent f2 f1 := by
  intro args res
  exact (h args res).symm

theorem semantic_equiv_trans {f1 f2 f3 : MirFunction}
    (h1 : SemanticEquivalent f1 f2) (h2 : SemanticEquivalent f2 f3) :
    SemanticEquivalent f1 f3 := by
  intro args res
  exact (h1 args res).trans (h2 args res)

theorem const_fold_stmt_equiv (env : MirEnv) (x y z : Local) (op : Op) (n1 n2 n3 : Int)
    (hy : lookup env y = some (Val.intVal n1))
    (hz : lookup env z = some (Val.intVal n2))
    (hop : evalOp op n1 n2 = some n3) :
    StmtStep env (Statement.Assign x (Rvalue.BinOp op y z)) (update env x (Val.intVal n3)) /\
    StmtStep env (Statement.Assign x (Rvalue.Constant n3)) (update env x (Val.intVal n3)) := by
  constructor
    -  exact StmtStep.assign_binop env x y z op n1 n2 n3 hy hz hop
    -  exact StmtStep.assign_const env x n3

theorem step_blocks_equiv {f1 f2 : MirFunction}
    (h_blocks : forall  b, b  in  f1.blocks <-> b  in  f2.blocks)
    {s s' : MirState} (h : Step f1 s s') : Step f2 s s' := by
  cases h with
  | stmt b st env' hb hpc hst hstmt =>
    exact Step.stmt s b st env' ((h_blocks b).mp hb) hpc hst hstmt
  | branch b target hb hpc hlen hterm =>
    exact Step.branch s b target ((h_blocks b).mp hb) hpc hlen hterm
  | branchIf_true b cond then_t else_t n hb hpc hlen hterm hcond hnz =>
    exact Step.branchIf_true s b cond then_t else_t n ((h_blocks b).mp hb) hpc hlen hterm hcond hnz
  | branchIf_false b cond then_t else_t hb hpc hlen hterm hcond =>
    exact Step.branchIf_false s b cond then_t else_t ((h_blocks b).mp hb) hpc hlen hterm hcond
  | switch b var targets default_t target hb hpc hlen hterm =>
    exact Step.switch s b var targets default_t target ((h_blocks b).mp hb) hpc hlen hterm
  | fork b left right join hb hpc hlen hterm =>
    exact Step.fork s b left right join ((h_blocks b).mp hb) hpc hlen hterm

theorem stepstar_blocks_equiv {f1 f2 : MirFunction}
    (h_blocks : forall  b, b  in  f1.blocks <-> b  in  f2.blocks)
    {s s' : MirState} (h : StepStar f1 s s') : StepStar f2 s s' := by
  induction h with
  | refl s => exact StepStar.refl s
  | step s1 s2 s3 hstep _ ih =>
    exact StepStar.step s1 s2 s3 (step_blocks_equiv h_blocks hstep) ih

theorem terminates_blocks_equiv {f1 f2 : MirFunction}
    (h_blocks : forall  b, b  in  f1.blocks <-> b  in  f2.blocks)
    {s : MirState} {res : Val} (h : TerminatesWith f1 s res) : TerminatesWith f2 s res := by
  cases h with
  | ret_some b retVar hb hpc hlen hterm hlookup =>
    exact TerminatesWith.ret_some b retVar ((h_blocks b).mp hb) hpc hlen hterm hlookup
  | ret_none b hb hpc hlen hterm hres =>
    exact TerminatesWith.ret_none b ((h_blocks b).mp hb) hpc hlen hterm hres

theorem semantic_equiv_of_blocks_equiv {f1 f2 : MirFunction}
    (h_entry : f1.entry = f2.entry)
    (h_blocks : forall  b, b  in  f1.blocks <-> b  in  f2.blocks) :
    SemanticEquivalent f1 f2 := by
  intro args res
  have h_init : initState f1 args = initState f2 args := by
    dsimp [initState]
    rw [h_entry]
  have h_blocks_rev : forall  b, b  in  f2.blocks <-> b  in  f1.blocks := by
    intro b
    exact (h_blocks b).symm
  constructor
    -  intro <s_f, h_star, h_term>
    rw [h_init] at h_star
    exact <s_f, stepstar_blocks_equiv h_blocks h_star, terminates_blocks_equiv h_blocks h_term>
    -  intro <s_f, h_star, h_term>
    rw [<- h_init] at h_star
    exact <s_f, stepstar_blocks_equiv h_blocks_rev h_star, terminates_blocks_equiv h_blocks_rev h_term>

end Supercompiler

\end{lstlisting}
